\subsection{Faithful and Irreducible Representations}\label{subsec:faithful-and-irreducible-representations}
Not all unital C*-algebras admit a faithful, irreducible representation, i.e. are primitive.

For example, unital, Abelian C*-algebras are ``rarely'' primitive. The following lemma formalizes this (See page 158 of Section 5.4 of \href{https://doi.org/10.1016/C2009-0-22289-6}{Murphy})
\begin{lemma}[Unital, Abelian C*-algebras are Rarely Primitive]
A non-zero Abelian C*-algebra $\mathcal{A}$ is primitive if and only if $\mathcal{A} = \mathbb{C} \mathbf{1}$.
\end{lemma}

So, primitive, Abelian C*-algebras are in this sense much more ``rare'' than Abelian C*-algebras.

As to the physical motivation in requiring $\mathfrak{U}$ be primitive, the original reference (\href{https://doi.org/10.1063/1.1704187}{Haag and Kastler}) doesn't provide much detail.

This original reference says the requirement is ``natural'' (Page 852 of \href{https://doi.org/10.1063/1.1704187}{Haag and Kastler}), and doesn't go much further in motivating the requirement.

The only additional motivation is on the same page in footnote 27 which references Appendix III of an unpublished article by Misra. Misra's article is claimed to exhibit a nonsimple C*-algebra that is of ``physical interest''. This is of relevance as simple C*-algebras are primitive (See page 158 of Section 5.4 of \href{https://doi.org/10.1016/C2009-0-22289-6}{Murphy}). So, one might think Axiom 6 should require $\mathfrak{U}$ be simple and not primitive. Misra's example is the counterweight to such an impulse.

All of this said, the physical motivation for the requirement of being primitive is lacking. In addition, in later versions of the axioms presented by Haag (\href{https://doi.org/10.1007/978-3-642-61458-3}{Haag}), this axiom is dropped completely.

We will follow suit and simply drop this axiom.
